% source: manti.01.ma, 1ère SM, série suites (Drive 0B6pfSJcEVng0MUFxTlR3Ml9vM2c)
% Ex. 1 = série française suites_s1 (non repris)
\ex{Suites}{MUFx}{Ex. 2}{9}
On considère la suite $(U_n)_{n\in\mathbb{N}}$ définie par $U_0=1$ et $U_{n+1}=\dfrac{6U_n-2}{2U_n+1}$.
\begin{enumerate}
\item Montrer que $(\forall n\in\mathbb{N})\ \dfrac12<U_n<2$.
\item Étudier la monotonie de la suite $(U_n)_{n\in\mathbb{N}}$.
\item On pose $V_n=\dfrac{2U_n-1}{U_n-2}$.
\begin{enumerate}
\item Montrer que $(V_n)_{n\geq0}$ est une suite géométrique.
\item Montrer que $U_n=\dfrac{2^n+2\times5^n}{2^{n+1}+5^n}$.
\end{enumerate}
\end{enumerate}
\fin
\ex{Suites}{MUFx}{Ex. 3}{11}
On considère la fonction $f$ définie par $f(x)=\dfrac{2x+3}{x+2}$.
\begin{enumerate}
\item \begin{enumerate}
\item Étudier la monotonie de $f$ sur l'intervalle $I=[1,2]$.
\item En déduire que $f(I)\subset I$.
\end{enumerate}
\item Soit $(U_n)_{n\in\mathbb{N}}$ la suite définie par $U_0=1$ et $U_{n+1}=f(U_n)$.
\begin{enumerate}
\item Calculer $U_1$ et montrer que $(\forall n\in\mathbb{N})\ 1\leq U_n<2$.
\item Montrer que la suite $(U_n)_{n\in\mathbb{N}}$ est croissante.
\item Montrer que $(\forall n\in\mathbb{N})\ \left|U_{n+1}-\sqrt3\right|\leq\dfrac12\left|U_n-\sqrt3\right|$.
\item Montrer par récurrence que $(\forall n\in\mathbb{N})\ \left|U_n-\sqrt3\right|\leq\left(\dfrac12\right)^n$.
\end{enumerate}
\end{enumerate}
\fin
\ex{Suites}{MUFx}{Ex. 4}{11}
\textbf{(I)} On considère la fonction $f$ définie sur $\mathbb{R}$ par $f(x)=\dfrac{x^2+x}{x^2+1}$ et on pose $I=[0,1]$.
\begin{enumerate}
\item Montrer que $f(I)\subset I$.
\item Montrer que $(\forall x\in I)\ f(x)\geq x$.
\end{enumerate}
\textbf{(II)} Soit $(U_n)_{n\in\mathbb{N}}$ la suite définie par $U_0=\dfrac12$ et $U_{n+1}=f(U_n)$.
\begin{enumerate}
\item Montrer que $(\forall n\in\mathbb{N})\ 0<U_n<1$.
\item Montrer que la suite $(U_n)_{n\in\mathbb{N}}$ est croissante.
\item \begin{enumerate}
\item Montrer que $(\forall n\in\mathbb{N})\ |U_{n+1}-1|\leq\dfrac45|U_n-1|$.
\item Montrer que $(\forall n\in\mathbb{N})\ |U_n-1|\leq\dfrac12\left(\dfrac45\right)^n$.
\end{enumerate}
\end{enumerate}
\fin
\ex{Suites}{MUFx}{Ex. 5}{12}
Soient $(U_n)_n$ et $(V_n)_n$ les suites définies par
$U_0=1,\ U_{n+1}=\dfrac{3U_n+5V_n}{8}$ \quad et \quad $V_0=2,\ V_{n+1}=\dfrac{U_n+5V_n}{6}$.
\begin{enumerate}
\item Calculer $U_1$ et $V_1$.
\item On pose $W_n=V_n-U_n$.
\begin{enumerate}
\item Montrer que $(W_n)_n$ est une suite géométrique.
\item Calculer la somme $S_n=W_0+W_1+\dots+W_{n-1}$.
\end{enumerate}
\item \begin{enumerate}
\item Vérifier que $(\forall n\in\mathbb{N}^*)\ U_n-U_{n-1}=\dfrac58W_{n-1}$.
\item En déduire que $U_n=U_0+\dfrac58S_n$, puis montrer que pour tout $n\in\mathbb{N}^*$ : $\left|U_n-\dfrac{34}{19}\right|\leq\left(\dfrac{5}{24}\right)^n$.
\end{enumerate}
\end{enumerate}
\fin
\ex{Suites}{MUFx}{Ex. 6}{12}
Soit $(U_n)_{n\in\mathbb{N}}$ la suite définie par $U_0=0$ et $U_{n+1}=\left(1+\dfrac12\sqrt{1+U_n}\right)^2-1$.
\begin{enumerate}
\item Montrer que $(\forall n\in\mathbb{N})\ 0\leq U_n<3$.
\item \begin{enumerate}
\item Montrer que $U_{n+1}-U_n=-\dfrac34\left[\left(\sqrt{1+U_n}-\dfrac23\right)^2-\dfrac{16}{9}\right]$.
\item Étudier la monotonie de la suite $(U_n)_{n\in\mathbb{N}}$.
\end{enumerate}
\item On pose $V_n=\sqrt{1+U_n}-2$.
\begin{enumerate}
\item Montrer que $(V_n)_n$ est une suite géométrique en précisant sa raison.
\item Calculer $U_n$ en fonction de $n$.
\item Calculer la somme $S_n=\displaystyle\sum_{k=0}^{n}\sqrt{1+U_k}$.
\end{enumerate}
\end{enumerate}
\fin
